\lampiransection{Bukti Observasi GrabFood Nasi Gerilya 13 Juni 2026}
\label{app:evidence-observasi-grabfood-20260613}

Lampiran ini memuat bukti visual dari observasi halaman GrabFood Nasi Gerilya - Glugur Kota pada 13 Juni 2026. Bukti visual digunakan untuk memperjelas temuan mengenai profil toko, rating, promo, menu, harga, ketersediaan produk, dan ulasan pelanggan yang dibahas pada Bab IV. Kode data dicantumkan agar hubungan antara bukti visual, ringkasan temuan, dan pembahasan dapat ditelusuri secara jelas.

Bukti visual dalam lampiran ini disajikan bersama keterangan singkat untuk mendukung proses pemeriksaan data dan penarikan temuan penelitian.

\begin{table}[!htbp]
    \centering
    \caption{Identitas Bukti Observasi GrabFood}
    \label{tab:identitas-evidence-observasi-grabfood-20260613}
    \begin{tabularx}{\textwidth}{|Z{3.4cm}|Y|}
        \toprule
        \textbf{Aspek} & \textbf{Keterangan} \\ \hline
        \midrule
        Objek observasi & Nasi Gerilya - Glugur Kota \\ \hline
        Platform & GrabFood \\ \hline
        Tanggal observasi & 13 Juni 2026 \\ \hline
        Fokus bukti & Profil toko, rating, jumlah penilaian, promo/voucher, menu, harga, foto produk, menu tidak tersedia, dan ulasan pelanggan \\ \hline
        Jumlah bukti visual & 76 bukti visual \\ \hline
        \bottomrule
    \end{tabularx}
    \tablesource{Diolah peneliti dari observasi GrabFood Nasi Gerilya tanggal 13 Juni 2026.}
\end{table}

\begin{table}[!htbp]
    \centering
    \caption{Ringkasan Kelompok Bukti Observasi}
    \label{tab:ringkasan-kelompok-evidence-grabfood-20260613}
    \footnotesize
    \setlength{\tabcolsep}{4pt}
    \renewcommand{\arraystretch}{0.95}
    \begin{tabularx}{\textwidth}{|Z{3cm}|c|Y|Y|}
        \toprule
        \textbf{Kelompok Bukti} & \textbf{Jumlah} & \textbf{Isi Data} & \textbf{Kegunaan Analisis} \\ \hline
        \midrule
        Profil dan rating & 4 & Nama toko, kategori, status toko, jarak, rating, jumlah penilaian, dan ringkasan ulasan agregat & Identitas objek, reputasi platform, dan bukti fisik digital \\ \hline
        Promo dan voucher & 19 & Promo persentase, potongan nominal, diskon ongkir, bank, dan \textit{group order} & Analisis promosi, insentif harga, dan peluang platform \\ \hline
        Menu, harga, dan foto & 24 & Menu utama, menu terlaris, menu promo, minuman, sambal, tambahan, harga, dan foto produk & Analisis produk, harga, variasi menu, dan tampilan digital \\ \hline
        Menu tidak tersedia & 6 & Menu atau produk yang terlihat berstatus \textit{Nggak tersedia} & Analisis ketersediaan stok, proses operasional, dan risiko pengalaman pelanggan \\ \hline
        Ringkasan harga & 3 & Rentang harga keseluruhan, rentang harga item tersedia, dan harga menu utama tersedia & Analisis harga dan pembeda harga aktif dibanding item tidak tersedia \\ \hline
        Ulasan positif & 10 & Ulasan pelanggan positif mengenai rasa, porsi, kesegaran, kebersihan, request, kemasan, dan pembelian rutin & Dasar identifikasi kekuatan internal dan loyalitas pelanggan \\ \hline
        Ulasan negatif/netral & 10 & Ulasan mengenai salah item, item kurang lengkap, add-on tidak diberikan, harga/porsi, kualitas, koordinasi staf, dan pengalaman pascakonsumsi & Dasar identifikasi kelemahan, risiko layanan, dan kebutuhan triangulasi \\ \hline
        \bottomrule
    \end{tabularx}
    \tablesource{Diolah peneliti dari observasi GrabFood tanggal 13 Juni 2026.}
\end{table}

\begin{table}[!htbp]
    \centering
    \caption{Pemetaan Kode Temuan dan Kode Data Pendukung Observasi GrabFood}
    \label{tab:temuan-utama-evidence-observasi-grabfood-20260613}
    \scriptsize
    \setlength{\tabcolsep}{3pt}
    \renewcommand{\arraystretch}{0.95}
    \begin{tabularx}{\textwidth}{|Z{1.55cm}|Z{2.45cm}|Y|Y|}
        \toprule
        \textbf{Kode Temuan} & \textbf{Temuan} & \textbf{Keterangan} & \textbf{Kode Data Pendukung} \\ \hline
        \midrule
        TO-GF-01 & Identitas dan reputasi toko & Nama toko, status toko, jarak, dan rating 4,7 dari 4.584 penilaian terlihat pada halaman GrabFood Nasi Gerilya. & GF-PROFILE-001; GF-PROFILE-002; GF-RATING-001 \\ \hline
        TO-GF-02 & Ringkasan ulasan agregat & Ringkasan ulasan agregat menyebut rasa enak, porsi besar, bahan segar, kemasan rapi, menu beragam, dan harga terjangkau. & GF-REVIEW-SUMMARY-001 \\ \hline
        TO-GF-03 & Promo aktif & Terdapat 19 promo/voucher yang terlihat pada halaman Promo GrabFood, termasuk diskon persentase, potongan nominal, diskon ongkir, bank, dan \textit{group order}. & GF-PROMO-001--GF-PROMO-019 \\ \hline
        TO-GF-04 & Menu terlaris & Menu berlabel terlaris yang terlihat adalah Nasi Padang Rendang Sapi, Nasi Padang Ayam Goreng Bumbu, dan Nasi Padang Ayam Rendang. & GF-MENU-TERLARIS-001--GF-MENU-TERLARIS-003 \\ \hline
        TO-GF-05 & Rentang harga aktif & Rentang harga item tersedia yang dikodekan adalah Rp1.100--Rp72.000, sedangkan paket nasi Padang tersedia berada pada Rp27.500--Rp72.000. & GF-PRICE-RANGE-001; GF-PRICE-RANGE-003 \\ \hline
        TO-GF-06 & Menu tidak tersedia & Terdapat 6 menu atau produk yang terlihat berstatus \textit{Nggak tersedia}, termasuk menu premium kepala ikan dan produk sambal kemasan. & GF-SOLDOUT-001--GF-SOLDOUT-006 \\ \hline
        TO-GF-07 & Tema ulasan positif & Pelanggan menonjolkan rasa, porsi, kesegaran bahan, kebersihan, pengemasan, request yang dipenuhi, dan kebiasaan membeli ulang. & GF-REV-POS-001--GF-REV-POS-010 \\ \hline
        TO-GF-08 & Tema ulasan negatif/netral & Keluhan berkaitan dengan salah item, item kurang lengkap, add-on tidak diberikan, persepsi harga/porsi, kesegaran, koordinasi staf, dan risiko pengalaman pascakonsumsi. & GF-REV-NEG-001--GF-REV-NEG-010 \\ \hline
        \bottomrule
    \end{tabularx}
    \tablesource{Diolah peneliti dari observasi GrabFood tanggal 13 Juni 2026.}
\end{table}

\clearpage
\subsection*{Bukti Visual Observasi GrabFood}

Bukti visual berikut disajikan sebagai pendukung temuan observasi GrabFood. Keterangan pada caption gambar menunjukkan kode data dan data utama yang digunakan pada ringkasan dan pembahasan Bab IV.

\newlength{\gfEvidenceMaxHeight}
\newcommand{\gfEvidenceBlock}[4]{%
    \setlength{\gfEvidenceMaxHeight}{#4\textheight}%
    \setlength{\gfEvidenceMaxHeight}{0.7\gfEvidenceMaxHeight}%
    \par\noindent\begin{minipage}{\textwidth}%
        \centering
        \includegraphics[width=0.62\textwidth,height=\gfEvidenceMaxHeight,keepaspectratio]{#3}\par
        \figuresource[0.62\textwidth]{Dokumentasi observasi GrabFood Nasi Gerilya tanggal 13 Juni 2026; diolah peneliti (2026).}
        \captionof{figure}{#1 -- #2}%
    \end{minipage}\par\vspace{7pt}%
}

\newcommand{\gfEvidenceCompactBlock}[4]{%
    \setlength{\gfEvidenceMaxHeight}{#4\textheight}%
    \setlength{\gfEvidenceMaxHeight}{0.58\gfEvidenceMaxHeight}%
    \par\noindent\begin{minipage}{\textwidth}%
        \centering
        \includegraphics[width=0.56\textwidth,height=\gfEvidenceMaxHeight,keepaspectratio]{#3}\par
        \figuresource[0.56\textwidth]{Dokumentasi observasi GrabFood Nasi Gerilya tanggal 13 Juni 2026; diolah peneliti (2026).}
        \captionof{figure}{#1 -- #2}%
    \end{minipage}\par\vspace{4pt}%
}

% File ini digenerate dari O-GF-20260613-coded-manifest.csv.
% Jangan edit manual jika manifest berubah; regenerate dari manifest.
\subsection*{Profil dan Rating}

\gfEvidenceBlock{GF-PROFILE-001}{Nasi Gerilya - Glugur Kota}{resource/05-fieldwork/02-observasi-grabfood/20260613-nasi-gerilya/coded-data/profile-rating/gf-profile-001-nasi-gerilya-glugur-kota.jpeg}{0.26}

\gfEvidenceBlock{GF-PROFILE-002}{Tutup; buka besok; jarak 6.6 km; badge Grab bintang 5 dan centang}{resource/05-fieldwork/02-observasi-grabfood/20260613-nasi-gerilya/coded-data/profile-rating/gf-profile-002-tutup-buka-besok-jarak-6-6-km-badge-grab-bintang-5-dan-centang.jpeg}{0.26}

\gfEvidenceBlock{GF-RATING-001}{Rating 4.7 dari 4.584 penilaian}{resource/05-fieldwork/02-observasi-grabfood/20260613-nasi-gerilya/coded-data/profile-rating/gf-rating-001-rating-4-7-dari-4-584-penilaian.jpeg}{0.26}

\gfEvidenceBlock{GF-REVIEW-SUMMARY-001}{AI summary: rasa enak, porsi besar, bahan segar, kemasan rapi, menu beragam, harga terjangkau}{resource/05-fieldwork/02-observasi-grabfood/20260613-nasi-gerilya/coded-data/profile-rating/gf-review-summary-001-ai-summary-rasa-enak-porsi-besar-bahan-segar-kemasan-rapi-menu-berag.jpeg}{0.40}

\clearpage
\subsection*{Promo dan Voucher}

\gfEvidenceBlock{GF-PROMO-001}{Diskon 17\% hingga Rp50.000; tanpa belanja minimum}{resource/05-fieldwork/02-observasi-grabfood/20260613-nasi-gerilya/coded-data/promo-voucher/gf-promo-001-diskon-17-hingga-rp50-000-tanpa-belanja-minimum.jpeg}{0.26}

\gfEvidenceBlock{GF-PROMO-002}{Diskon Rp15.000 BSI; minimum belanja terlihat terpotong}{resource/05-fieldwork/02-observasi-grabfood/20260613-nasi-gerilya/coded-data/promo-voucher/gf-promo-002-diskon-rp15-000-bsi-minimum-belanja-terlihat-terpotong.jpeg}{0.26}

\gfEvidenceBlock{GF-PROMO-003}{Diskon Rp10.000 BTN; minimum belanja terlihat terpotong}{resource/05-fieldwork/02-observasi-grabfood/20260613-nasi-gerilya/coded-data/promo-voucher/gf-promo-003-diskon-rp10-000-btn-minimum-belanja-terlihat-terpotong.jpeg}{0.26}

\gfEvidenceBlock{GF-PROMO-004}{Diskon Rp25.000 Bank Mega Syariah; minimum belanja terlihat terpotong}{resource/05-fieldwork/02-observasi-grabfood/20260613-nasi-gerilya/coded-data/promo-voucher/gf-promo-004-diskon-rp25-000-bank-mega-syariah-minimum-belanja-terlihat-terpotong.jpeg}{0.26}

\gfEvidenceBlock{GF-PROMO-005}{Diskon Rp4.500 untuk Gulai Cincang Gerilya}{resource/05-fieldwork/02-observasi-grabfood/20260613-nasi-gerilya/coded-data/promo-voucher/gf-promo-005-diskon-rp4-500-untuk-gulai-cincang-gerilya.jpeg}{0.26}

\gfEvidenceBlock{GF-PROMO-006}{Diskon Rp15.000 Bank Mega; minimum belanja terlihat terpotong}{resource/05-fieldwork/02-observasi-grabfood/20260613-nasi-gerilya/coded-data/promo-voucher/gf-promo-006-diskon-rp15-000-bank-mega-minimum-belanja-terlihat-terpotong.jpeg}{0.26}

\gfEvidenceBlock{GF-PROMO-007}{Diskon Rp25.000 Bank Mega; minimum belanja terlihat terpotong}{resource/05-fieldwork/02-observasi-grabfood/20260613-nasi-gerilya/coded-data/promo-voucher/gf-promo-007-diskon-rp25-000-bank-mega-minimum-belanja-terlihat-terpotong.jpeg}{0.26}

\gfEvidenceBlock{GF-PROMO-008}{Diskon 5\% untuk Group Order; minimum orang bergabung 3 ke 6}{resource/05-fieldwork/02-observasi-grabfood/20260613-nasi-gerilya/coded-data/promo-voucher/gf-promo-008-diskon-5-untuk-group-order-minimum-orang-bergabung-3-ke-6.jpeg}{0.26}

\gfEvidenceBlock{GF-PROMO-009}{Diskon Rp4.000 Bank Mega; minimum belanja terlihat terpotong}{resource/05-fieldwork/02-observasi-grabfood/20260613-nasi-gerilya/coded-data/promo-voucher/gf-promo-009-diskon-rp4-000-bank-mega-minimum-belanja-terlihat-terpotong.jpeg}{0.26}

\gfEvidenceBlock{GF-PROMO-010}{Diskon Rp10.000 Bank Mega; minimum belanja terlihat terpotong}{resource/05-fieldwork/02-observasi-grabfood/20260613-nasi-gerilya/coded-data/promo-voucher/gf-promo-010-diskon-rp10-000-bank-mega-minimum-belanja-terlihat-terpotong.jpeg}{0.26}

\gfEvidenceBlock{GF-PROMO-011}{Diskon ongkir Rp3.000; tanpa belanja minimum}{resource/05-fieldwork/02-observasi-grabfood/20260613-nasi-gerilya/coded-data/promo-voucher/gf-promo-011-diskon-ongkir-rp3-000-tanpa-belanja-minimum.jpeg}{0.26}

\gfEvidenceBlock{GF-PROMO-012}{Diskon Rp25.000 BSI; minimum belanja terlihat terpotong}{resource/05-fieldwork/02-observasi-grabfood/20260613-nasi-gerilya/coded-data/promo-voucher/gf-promo-012-diskon-rp25-000-bsi-minimum-belanja-terlihat-terpotong.jpeg}{0.26}

\gfEvidenceBlock{GF-PROMO-013}{Diskon Rp25.000 Raya Bank; minimum belanja terlihat terpotong}{resource/05-fieldwork/02-observasi-grabfood/20260613-nasi-gerilya/coded-data/promo-voucher/gf-promo-013-diskon-rp25-000-raya-bank-minimum-belanja-terlihat-terpotong.jpeg}{0.26}

\gfEvidenceBlock{GF-PROMO-014}{Diskon 25\% s.d. Rp100.000; minimum belanja Rp30.000}{resource/05-fieldwork/02-observasi-grabfood/20260613-nasi-gerilya/coded-data/promo-voucher/gf-promo-014-diskon-25-s-d-rp100-000-minimum-belanja-rp30-000.jpeg}{0.26}

\gfEvidenceBlock{GF-PROMO-015}{Diskon Rp15.000 BTN; minimum belanja terlihat terpotong}{resource/05-fieldwork/02-observasi-grabfood/20260613-nasi-gerilya/coded-data/promo-voucher/gf-promo-015-diskon-rp15-000-btn-minimum-belanja-terlihat-terpotong.jpeg}{0.26}

\gfEvidenceBlock{GF-PROMO-016}{Diskon ongkir Rp8.000; minimum belanja terlihat terpotong}{resource/05-fieldwork/02-observasi-grabfood/20260613-nasi-gerilya/coded-data/promo-voucher/gf-promo-016-diskon-ongkir-rp8-000-minimum-belanja-terlihat-terpotong.jpeg}{0.26}

\gfEvidenceBlock{GF-PROMO-017}{Diskon ongkir Rp8.000; minimum belanja terlihat terpotong}{resource/05-fieldwork/02-observasi-grabfood/20260613-nasi-gerilya/coded-data/promo-voucher/gf-promo-017-diskon-ongkir-rp8-000-minimum-belanja-terlihat-terpotong.jpeg}{0.26}

\gfEvidenceBlock{GF-PROMO-018}{Diskon Rp20.000 SBI card; minimum belanja terlihat terpotong}{resource/05-fieldwork/02-observasi-grabfood/20260613-nasi-gerilya/coded-data/promo-voucher/gf-promo-018-diskon-rp20-000-sbi-card-minimum-belanja-terlihat-terpotong.jpeg}{0.26}

\gfEvidenceBlock{GF-PROMO-019}{Diskon Rp20.000 Bank Jakarta; minimum belanja terlihat terpotong}{resource/05-fieldwork/02-observasi-grabfood/20260613-nasi-gerilya/coded-data/promo-voucher/gf-promo-019-diskon-rp20-000-bank-jakarta-minimum-belanja-terlihat-terpotong.jpeg}{0.26}

\clearpage
\subsection*{Menu, Harga, dan Foto Produk}

\gfEvidenceBlock{GF-MENU-TERATAS-001}{Nasi Padang Bungkus Gulai Cincang; Rp50.500 dari Rp55.000}{resource/05-fieldwork/02-observasi-grabfood/20260613-nasi-gerilya/coded-data/menu-harga-foto/gf-menu-teratas-001-nasi-padang-bungkus-gulai-cincang-rp50-500-dari-rp55-000.jpeg}{0.26}

\gfEvidenceBlock{GF-MENU-TERATAS-002}{Nasi Padang Kari Kambing GERILYA; Rp72.000}{resource/05-fieldwork/02-observasi-grabfood/20260613-nasi-gerilya/coded-data/menu-harga-foto/gf-menu-teratas-002-nasi-padang-kari-kambing-gerilya-rp72-000.jpeg}{0.26}

\gfEvidenceBlock{GF-MENU-TERATAS-003}{Nasi Padang Dendeng Sapi Bakar GERILYA; Rp56.000}{resource/05-fieldwork/02-observasi-grabfood/20260613-nasi-gerilya/coded-data/menu-harga-foto/gf-menu-teratas-003-nasi-padang-dendeng-sapi-bakar-gerilya-rp56-000.jpeg}{0.26}

\gfEvidenceBlock{GF-MENU-TERLARIS-001}{Nasi Padang Rendang Sapi GERILYA; Rp56.000; label Terlaris}{resource/05-fieldwork/02-observasi-grabfood/20260613-nasi-gerilya/coded-data/menu-harga-foto/gf-menu-terlaris-001-nasi-padang-rendang-sapi-gerilya-rp56-000-label-terlaris.jpeg}{0.26}

\gfEvidenceBlock{GF-MENU-TERLARIS-002}{Nasi Padang Ayam Goreng Bumbu GERILYA; Rp45.000; label Terlaris}{resource/05-fieldwork/02-observasi-grabfood/20260613-nasi-gerilya/coded-data/menu-harga-foto/gf-menu-terlaris-002-nasi-padang-ayam-goreng-bumbu-gerilya-rp45-000-label-terlaris.jpeg}{0.26}

\gfEvidenceBlock{GF-MENU-TERLARIS-003}{Nasi Padang Ayam Rendang GERILYA; Rp46.000; label Terlaris}{resource/05-fieldwork/02-observasi-grabfood/20260613-nasi-gerilya/coded-data/menu-harga-foto/gf-menu-terlaris-003-nasi-padang-ayam-rendang-gerilya-rp46-000-label-terlaris.jpeg}{0.26}

\gfEvidenceBlock{GF-MENU-PROMO-002}{Nasi Padang Gulai Ayam Kalasan GERILYA; Rp41.500 dari Rp46.000}{resource/05-fieldwork/02-observasi-grabfood/20260613-nasi-gerilya/coded-data/menu-harga-foto/gf-menu-promo-002-nasi-padang-gulai-ayam-kalasan-gerilya-rp41-500-dari-rp46-000.jpeg}{0.26}

\gfEvidenceBlock{GF-MENU-UTAMA-001}{Nasi Padang Ayam Bakar GERILYA; Rp45.000}{resource/05-fieldwork/02-observasi-grabfood/20260613-nasi-gerilya/coded-data/menu-harga-foto/gf-menu-utama-001-nasi-padang-ayam-bakar-gerilya-rp45-000.jpeg}{0.26}

\gfEvidenceBlock{GF-MENU-UTAMA-002}{Nasi Padang Cumi Nagih GERILYA; Rp70.000}{resource/05-fieldwork/02-observasi-grabfood/20260613-nasi-gerilya/coded-data/menu-harga-foto/gf-menu-utama-002-nasi-padang-cumi-nagih-gerilya-rp70-000.jpeg}{0.26}

\gfEvidenceBlock{GF-MENU-UTAMA-003}{Nasi Padang Ikan Lele Goreng Panas GERILYA; Rp39.000}{resource/05-fieldwork/02-observasi-grabfood/20260613-nasi-gerilya/coded-data/menu-harga-foto/gf-menu-utama-003-nasi-padang-ikan-lele-goreng-panas-gerilya-rp39-000.jpeg}{0.26}

\gfEvidenceBlock{GF-MENU-UTAMA-004}{Nasi Padang Ikan Tongkol Tuna Goreng Panas GERILYA; Rp48.000}{resource/05-fieldwork/02-observasi-grabfood/20260613-nasi-gerilya/coded-data/menu-harga-foto/gf-menu-utama-004-nasi-padang-ikan-tongkol-tuna-goreng-panas-gerilya-rp48-000.jpeg}{0.26}

\gfEvidenceBlock{GF-MENU-UTAMA-005}{Nasi Padang Ikan Gembung Goreng GERILYA; Rp50.000}{resource/05-fieldwork/02-observasi-grabfood/20260613-nasi-gerilya/coded-data/menu-harga-foto/gf-menu-utama-005-nasi-padang-ikan-gembung-goreng-gerilya-rp50-000.jpeg}{0.26}

\gfEvidenceBlock{GF-MENU-UTAMA-006}{Nasi Padang Ikan Nila Goreng GERILYA; Rp43.000}{resource/05-fieldwork/02-observasi-grabfood/20260613-nasi-gerilya/coded-data/menu-harga-foto/gf-menu-utama-006-nasi-padang-ikan-nila-goreng-gerilya-rp43-000.jpeg}{0.26}

\gfEvidenceBlock{GF-MENU-UTAMA-007}{Nasi Padang Telur Dadar GERILYA; Rp32.000}{resource/05-fieldwork/02-observasi-grabfood/20260613-nasi-gerilya/coded-data/menu-harga-foto/gf-menu-utama-007-nasi-padang-telur-dadar-gerilya-rp32-000.jpeg}{0.26}

\gfEvidenceBlock{GF-MENU-UTAMA-008}{Nasi Padang Berkedel Kentang GERILYA; Rp37.400}{resource/05-fieldwork/02-observasi-grabfood/20260613-nasi-gerilya/coded-data/menu-harga-foto/gf-menu-utama-008-nasi-padang-berkedel-kentang-gerilya-rp37-400.jpeg}{0.26}

\gfEvidenceBlock{GF-MENU-UTAMA-009}{Nasi Padang Sayur Polos GERILYA; Rp27.500}{resource/05-fieldwork/02-observasi-grabfood/20260613-nasi-gerilya/coded-data/menu-harga-foto/gf-menu-utama-009-nasi-padang-sayur-polos-gerilya-rp27-500.jpeg}{0.26}

\gfEvidenceBlock{GF-MENU-MINUMAN-001}{Air Mineral 400ml; Rp13.200}{resource/05-fieldwork/02-observasi-grabfood/20260613-nasi-gerilya/coded-data/menu-harga-foto/gf-menu-minuman-001-air-mineral-400ml-rp13-200.jpeg}{0.26}

\gfEvidenceBlock{GF-MENU-MINUMAN-002}{Es Teh Tawar; Rp13.200}{resource/05-fieldwork/02-observasi-grabfood/20260613-nasi-gerilya/coded-data/menu-harga-foto/gf-menu-minuman-002-es-teh-tawar-rp13-200.jpeg}{0.26}

\gfEvidenceBlock{GF-MENU-MINUMAN-003}{Es Teh Manis; Rp15.950}{resource/05-fieldwork/02-observasi-grabfood/20260613-nasi-gerilya/coded-data/menu-harga-foto/gf-menu-minuman-003-es-teh-manis-rp15-950.jpeg}{0.26}

\gfEvidenceBlock{GF-MENU-MINUMAN-004}{Es Jeruk Limau; Rp22.000}{resource/05-fieldwork/02-observasi-grabfood/20260613-nasi-gerilya/coded-data/menu-harga-foto/gf-menu-minuman-004-es-jeruk-limau-rp22-000.jpeg}{0.26}

\gfEvidenceBlock{GF-MENU-TAMBAHAN-001}{Nambah KUAH CAMPUR; Rp2.500}{resource/05-fieldwork/02-observasi-grabfood/20260613-nasi-gerilya/coded-data/menu-harga-foto/gf-menu-tambahan-001-nambah-kuah-campur-rp2-500.jpeg}{0.26}

\gfEvidenceBlock{GF-MENU-TAMBAHAN-002}{Nambah Kuah Gulai Kuning; Rp1.100}{resource/05-fieldwork/02-observasi-grabfood/20260613-nasi-gerilya/coded-data/menu-harga-foto/gf-menu-tambahan-002-nambah-kuah-gulai-kuning-rp1-100.jpeg}{0.26}

\gfEvidenceBlock{GF-MENU-SAMBAL-001}{GERILYA Sambal Hijau; Rp17.500}{resource/05-fieldwork/02-observasi-grabfood/20260613-nasi-gerilya/coded-data/menu-harga-foto/gf-menu-sambal-001-gerilya-sambal-hijau-rp17-500.jpeg}{0.26}

\gfEvidenceBlock{GF-MENU-SAMBAL-002}{GERILYA Sambal Andaliman Getir Kemasan; Rp38.500}{resource/05-fieldwork/02-observasi-grabfood/20260613-nasi-gerilya/coded-data/menu-harga-foto/gf-menu-sambal-002-gerilya-sambal-andaliman-getir-kemasan-rp38-500.jpeg}{0.26}

\clearpage
\subsection*{Menu Tidak Tersedia}

\gfEvidenceBlock{GF-SOLDOUT-001}{Gulai Merah Kepala Ikan GERILYA L; Rp275.000; Nggak tersedia}{resource/05-fieldwork/02-observasi-grabfood/20260613-nasi-gerilya/coded-data/sold-out/gf-soldout-001-gulai-merah-kepala-ikan-gerilya-l-rp275-000-nggak-tersedia.jpeg}{0.32}

\gfEvidenceBlock{GF-SOLDOUT-002}{Gulai Merah Kepala Ikan GERILYA M; Rp248.050; Nggak tersedia}{resource/05-fieldwork/02-observasi-grabfood/20260613-nasi-gerilya/coded-data/sold-out/gf-soldout-002-gulai-merah-kepala-ikan-gerilya-m-rp248-050-nggak-tersedia.jpeg}{0.26}

\gfEvidenceBlock{GF-SOLDOUT-003}{Nasi Padang Gembung Kuring Bakar GERILYA; Rp47.850; Nggak tersedia}{resource/05-fieldwork/02-observasi-grabfood/20260613-nasi-gerilya/coded-data/sold-out/gf-soldout-003-nasi-padang-gembung-kuring-bakar-gerilya-rp47-850-nggak-tersedia.jpeg}{0.26}

\gfEvidenceBlock{GF-SOLDOUT-004}{Nasi Padang Ayam Pop Sauce Creamy GERILYA; Rp46.000; Nggak tersedia}{resource/05-fieldwork/02-observasi-grabfood/20260613-nasi-gerilya/coded-data/sold-out/gf-soldout-004-nasi-padang-ayam-pop-sauce-creamy-gerilya-rp46-000-nggak-tersedia.jpeg}{0.26}

\gfEvidenceBlock{GF-SOLDOUT-005}{GERILYA Sambal Cumi Nagih Kemasan; Rp38.500; Nggak tersedia}{resource/05-fieldwork/02-observasi-grabfood/20260613-nasi-gerilya/coded-data/sold-out/gf-soldout-005-gerilya-sambal-cumi-nagih-kemasan-rp38-500-nggak-tersedia.jpeg}{0.26}

\gfEvidenceBlock{GF-SOLDOUT-006}{GERILYA Sambal Tuna Asap Keumamah Aceh Kemasan; Rp38.500; Nggak tersedia}{resource/05-fieldwork/02-observasi-grabfood/20260613-nasi-gerilya/coded-data/sold-out/gf-soldout-006-gerilya-sambal-tuna-asap-keumamah-aceh-kemasan-rp38-500-nggak-tersedia.jpeg}{0.26}

\clearpage
\subsection*{Ringkasan Harga}

\gfEvidenceBlock{GF-PRICE-RANGE-001}{Rentang harga observasi: Rp1.100 sampai Rp275.000; tersedia: Rp1.100 sampai Rp72.000; paket nasi padang tersedia: Rp27.500 sampai Rp72.000}{resource/05-fieldwork/02-observasi-grabfood/20260613-nasi-gerilya/coded-data/ringkasan-harga/gf-price-range-001-rentang-harga-observasi-rp1-100-sampai-rp275-000-tersedia-rp1-100-sampa.jpeg}{0.40}

\gfEvidenceBlock{GF-PRICE-RANGE-002}{Harga tertinggi terlihat: Gulai Merah Kepala Ikan L Rp275.000 dan M Rp248.050, tetapi tidak tersedia}{resource/05-fieldwork/02-observasi-grabfood/20260613-nasi-gerilya/coded-data/ringkasan-harga/gf-price-range-002-harga-tertinggi-terlihat-gulai-merah-kepala-ikan-l-rp275-000-dan-m-rp24.jpeg}{0.40}

\gfEvidenceBlock{GF-PRICE-RANGE-003}{Harga atas menu utama tersedia: Kari Kambing Rp72.000 dan Dendeng Sapi Bakar Rp56.000}{resource/05-fieldwork/02-observasi-grabfood/20260613-nasi-gerilya/coded-data/ringkasan-harga/gf-price-range-003-harga-atas-menu-utama-tersedia-kari-kambing-rp72-000-dan-dendeng-sapi-b.jpeg}{0.40}

\clearpage
\subsection*{Ulasan Positif}

\gfEvidenceBlock{GF-REV-POS-001}{Muslim M.: keren ni, konsisten rasa}{resource/05-fieldwork/02-observasi-grabfood/20260613-nasi-gerilya/coded-data/ulasan-positif/gf-rev-pos-001-muslim-m-keren-ni-konsisten-rasa.jpeg}{0.32}

\gfEvidenceBlock{GF-REV-POS-002}{Raymond: Ayam popnya wenak dan creamy}{resource/05-fieldwork/02-observasi-grabfood/20260613-nasi-gerilya/coded-data/ulasan-positif/gf-rev-pos-002-raymond-ayam-popnya-wenak-dan-creamy.jpeg}{0.32}

\gfEvidenceBlock{GF-REV-POS-003}{tasyaa: rasanya enak, sayur fresh, porsi agak kebanyakan}{resource/05-fieldwork/02-observasi-grabfood/20260613-nasi-gerilya/coded-data/ulasan-positif/gf-rev-pos-003-tasyaa-rasanya-enak-sayur-fresh-porsi-agak-kebanyakan.jpeg}{0.50}

\gfEvidenceBlock{GF-REV-POS-004}{Wie W.: langganan keluarga; rasa enak, porsi besar, bersih, request dipenuhi}{resource/05-fieldwork/02-observasi-grabfood/20260613-nasi-gerilya/coded-data/ulasan-positif/gf-rev-pos-004-wie-w-langganan-keluarga-rasa-enak-porsi-besar-bersih-request-dipenuhi.jpeg}{0.40}

\gfEvidenceBlock{GF-REV-POS-005}{Wie W.: langgananku setiap Minggu, Best}{resource/05-fieldwork/02-observasi-grabfood/20260613-nasi-gerilya/coded-data/ulasan-positif/gf-rev-pos-005-wie-w-langgananku-setiap-minggu-best.jpeg}{0.32}

\gfEvidenceBlock{GF-REV-POS-006}{ita: ikan gembung dan lele fresh; rasanya enak}{resource/05-fieldwork/02-observasi-grabfood/20260613-nasi-gerilya/coded-data/ulasan-positif/gf-rev-pos-006-ita-ikan-gembung-dan-lele-fresh-rasanya-enak.jpeg}{0.32}

\gfEvidenceBlock{GF-REV-POS-007}{Lina: semua menu enak, semua recommended}{resource/05-fieldwork/02-observasi-grabfood/20260613-nasi-gerilya/coded-data/ulasan-positif/gf-rev-pos-007-lina-semua-menu-enak-semua-recommended.jpeg}{0.32}

\gfEvidenceBlock{GF-REV-POS-008}{Gusti: best karena sayur dan lauk dipisah}{resource/05-fieldwork/02-observasi-grabfood/20260613-nasi-gerilya/coded-data/ulasan-positif/gf-rev-pos-008-gusti-best-karena-sayur-dan-lauk-dipisah.jpeg}{0.32}

\gfEvidenceBlock{GF-REV-POS-009}{Derrick: rasanya enak, porsinya besar}{resource/05-fieldwork/02-observasi-grabfood/20260613-nasi-gerilya/coded-data/ulasan-positif/gf-rev-pos-009-derrick-rasanya-enak-porsinya-besar.jpeg}{0.32}

\gfEvidenceBlock{GF-REV-POS-010}{Suba: enak banget menurut keluarga}{resource/05-fieldwork/02-observasi-grabfood/20260613-nasi-gerilya/coded-data/ulasan-positif/gf-rev-pos-010-suba-enak-banget-menurut-keluarga.jpeg}{0.32}

\clearpage
\subsection*{Ulasan Negatif/Netral}

\gfEvidenceCompactBlock{GF-REV-NEG-001}{albert p.: Jangek siram tidak dikasih; restoran meminta maaf dan berjanji memperhatikan kelengkapan}{resource/05-fieldwork/02-observasi-grabfood/20260613-nasi-gerilya/coded-data/ulasan-negatif-netral/gf-rev-neg-001-albert-p-jangek-siram-tidak-dikasih-restoran-meminta-maaf-dan-berjanji-memp.jpeg}{0.62}

\gfEvidenceCompactBlock{GF-REV-NEG-002}{Talia: pesan ayam rendang, dikirim ayam goreng}{resource/05-fieldwork/02-observasi-grabfood/20260613-nasi-gerilya/coded-data/ulasan-negatif-netral/gf-rev-neg-002-talia-pesan-ayam-rendang-dikirim-ayam-goreng.jpeg}{0.32}

\gfEvidenceCompactBlock{GF-REV-NEG-003}{irawaty: udang goreng hanya 3 biji, terlalu mahal, tidak worth it}{resource/05-fieldwork/02-observasi-grabfood/20260613-nasi-gerilya/coded-data/ulasan-negatif-netral/gf-rev-neg-003-irawaty-udang-goreng-hanya-3-biji-terlalu-mahal-tidak-worth-it.jpeg}{0.32}

\gfEvidenceCompactBlock{GF-REV-NEG-004}{Hendra: pesan ikan kakap tetapi diberikan ikan nila; restoran meminta maaf}{resource/05-fieldwork/02-observasi-grabfood/20260613-nasi-gerilya/coded-data/ulasan-negatif-netral/gf-rev-neg-004-hendra-pesan-ikan-kakap-tetapi-diberikan-ikan-nila-restoran-meminta-maaf.jpeg}{0.62}

\gfEvidenceCompactBlock{GF-REV-NEG-005}{siyaga.: makanan enak tetapi kuah yang diminta/dibayar tidak diberikan; restoran merespons}{resource/05-fieldwork/02-observasi-grabfood/20260613-nasi-gerilya/coded-data/ulasan-negatif-netral/gf-rev-neg-005-siyaga-makanan-enak-tetapi-kuah-yang-diminta-dibayar-tidak-diberikan-restor.jpeg}{0.62}

\gfEvidenceCompactBlock{GF-REV-NEG-006}{lily: nasi agak keras, keripik kentang melempem, packaging bagus}{resource/05-fieldwork/02-observasi-grabfood/20260613-nasi-gerilya/coded-data/ulasan-negatif-netral/gf-rev-neg-006-lily-nasi-agak-keras-keripik-kentang-melempem-packaging-bagus.jpeg}{0.32}

\gfEvidenceCompactBlock{GF-REV-NEG-007}{Khash: not quite fresh; tough meat; okay but does not meet expectation}{resource/05-fieldwork/02-observasi-grabfood/20260613-nasi-gerilya/coded-data/ulasan-negatif-netral/gf-rev-neg-007-khash-not-quite-fresh-tough-meat-okay-but-does-not-meet-expectation.jpeg}{0.40}

\gfEvidenceCompactBlock{GF-REV-NEG-008}{Hartini: sayur nangka ada rada basi}{resource/05-fieldwork/02-observasi-grabfood/20260613-nasi-gerilya/coded-data/ulasan-negatif-netral/gf-rev-neg-008-hartini-sayur-nangka-ada-rada-basi.jpeg}{0.32}

\gfEvidenceCompactBlock{GF-REV-NEG-009}{Vivi: staf kurang koordinasi; salah order ke driver; pelanggan keluar ongkos lagi; rasa biasa saja}{resource/05-fieldwork/02-observasi-grabfood/20260613-nasi-gerilya/coded-data/ulasan-negatif-netral/gf-rev-neg-009-vivi-staf-kurang-koordinasi-salah-order-ke-driver-pelanggan-keluar-ongkos-l.jpeg}{0.40}

\gfEvidenceCompactBlock{GF-REV-NEG-010}{Fajar: setelah makan langsung diare; restoran memberi respons}{resource/05-fieldwork/02-observasi-grabfood/20260613-nasi-gerilya/coded-data/ulasan-negatif-netral/gf-rev-neg-010-fajar-setelah-makan-langsung-diare-restoran-memberi-respons.jpeg}{0.40}


